%%%%%%
%
% $Author: Abishek,Bruna,Fedor $
% $Date: 07.10.2024 $
% $Path: ML24-02-Gait-Pattern-Recognition\report
% $Version: 1.0 $
%
%
% !TeX encoding = utf8
% !TeX root = Rename
% !TeX TXS-program:bibliography = txs:///bibtex
%
%%%%%%

\chapter{Introduction to the project}

\section{General}

In recent years, wearable devices have become more popular among the general population as they allow for continuously tracking health parameters and even enable for early detection of disease \cite{Sreenilayam:2020}The gait pattern is a parameter that is used to evaluate one’s walking and in medical research has been used as an indicator of health. Gait pattern research has been conducted in various areas to aid in elderly care such in fall detection for patients with Parkinson’s disease per \cite{Perezsiguas:2023}.
Furthermore, one of the challenges elderly adults face is an increased risks of falls and health complication as aging progressively impairs the person’s gait (\cite{Perezsiguas:2023}). \cite{Perezsiguas:2023} suggests that the available fall prevention systems are expensive, fragile and complex to handle, creating a need for more practical, simple and cheap systems to evaluate and detect abnormal gait that causes falls in the elderly people with diseases that impair the gait.
This project aims to create a wearable device that detect abnormal gait patterns associated with Stroke and Parkinson’s disease to alert its user or caregiver of the gait abnormality which can point to an increased risk for falls; a need to consult a doctor; and/or the need to find devices that can assist with walking such as walkers.




\bigskip


\section{Challenges and Proposed Solutions}

\subsection{Challenges}

\textbf{Challenge 1: Sensor Noise and Drift}  
The IMU sensor in the Arduino Nano 33 BLE Sense is prone to noise and drift caused by environmental factors such as temperature changes, electrical interference, and mechanical imperfections \cite{Bosch:2024}. This can distort motion data, reducing the reliability of gait classification.

\textbf{Challenge 2: Real-Time Data Processing}  
The Arduino Nano 33 BLE Sense has limited computational power and memory, making it challenging to process ML models and large datasets in real time. These constraints necessitate optimization of the ML pipeline for resource efficiency\cite{David:2021}.  

\textbf{Challenge 3: Variability in Gait Patterns}  
Human gait is highly variable due to differences in walking speed, terrain, and individual physical conditions. This variability can lead to difficulties in distinguishing between normal and abnormal gait patterns without introducing false positives\cite{Hausdorff:2009}. 

\textbf{Challenge 4: Hardware and Device Design}  
Designing a lightweight, compact, and energy-efficient wearable device poses challenges. The system must incorporate a reliable power source, maintain functionality over extended periods, and ensure user comfort without compromising performance \cite{ArduinoNano33Sense2025}.

\textbf{Challenge 5: Absence of a Parkinson’s Disease Subject}  
The team did not have access to a person with Parkinson’s disease for data collection. This limitation made it challenging to gather realistic, disease-specific gait patterns for model training and validation.

\textbf{Challenge 6: Dataset Limitations}  
Most publicly available gait datasets use two or more sensors to collect motion data, whereas this project aimed to work with a single IMU sensor. Identifying a reliable dataset matching the project’s requirements was a significant challenge.

\subsection{Proposed Solutions}

\textbf{Solution 1: Sensor Noise and Drift}  
Dynamic filtering methods such as the Kalman Filter were implemented to reduce noise and improve the accuracy of motion data. These filters dynamically adjust parameters in real time to ensure precise and reliable sensor readings \cite{Kalman:2012}.

\textbf{Solution 2: Real-Time Data Processing}  
TinyML techniques were applied to optimize the ML models for deployment on resource-constrained hardware. By reducing computational complexity, the models were adapted to run efficiently on the Arduino Nano 33 BLE Sense while maintaining accuracy \cite{David:2021, Abbas:2023}.

\textbf{Solution 3: Variability in Gait Patterns}  
The model was trained using TensorFlow, which provides flexible and scalable capabilities for handling diverse gait patterns. Advanced model architectures and abstraction layers allowed the system to accommodate the variability inherent in human gait \cite{Kumari2024}.

\textbf{Solution 4: Hardware and Device Design}  
The hardware design included a compact and lightweight form factor powered by 6V battery and embedded DC-DC converter. This configuration ensured optimal energy efficiency and long-term usability \cite{Tarduino:2020}.

\textbf{Solution 5: Absence of a Parkinson’s Disease Subject}  
A publicly available dataset containing gait patterns specific to Parkinson’s disease was utilized to overcome the lack of access to a subject with the condition. This dataset provided sufficient data for training and evaluating the ML model \cite{ZhangWenwen:2022}.

\textbf{Solution 6: Dataset Limitations}  
A suitable dataset utilizing a single IMU sensor was identified to align with the project’s hardware setup. The dataset included relevant gait parameters such as acceleration, step length, and stride time, ensuring compatibility with the project’s objectives \cite{Panebianco:2018}.

\section{Structure of the Report}

The structure of this report has been designed to provide a comprehensive understanding of the challenges, methodologies, and outcomes related to the development of a gait pattern recognition system using the Arduino Nano 33 BLE Sense . The report begins with an introduction that highlights the project's objectives, the technical hurdles faced, and the proposed solutions. Challenges such as sensor noise, variability in gait data, and hardware constraints are addressed, with a focus on employing machine learning models deployed on embedded systems to overcome these issues \cite{Kumari2024}.

The report is divided into the following sections:

\subsection*{Introduction}
The first chapter introduces the project, outlining its purpose and the primary goals. This includes:
\begin{itemize}
    \item A general overview of gait pattern recognition and its applications.
    \item A detailed explanation of the report's structure.
    \item Identification of key challenges and proposed solutions, including hardware and software integration.
\end{itemize}

\subsection*{Domain Knowledge}
This section delves into the domain understanding required for gait classification, focusing on:
\begin{itemize}
    \item The application of machine learning in gait pattern recognition.
    \item Challenges in data acquisition, quality, and relevance.
    \item Identification of anomalies and outliers in the dataset.
\end{itemize}

\subsection*{Hardware Description}
A detailed analysis of the Arduino Nano 33 BLE Sense  is presented, covering:
\begin{itemize}
    \item Key features and technical specifications of the board.
    \item Descriptions of on-board sensors, including the IMU (LSM9DS1), accelerometer,gyroscope and magnetometer.
    \item A comparison of hardware revisions (Rev1, Rev2, and Lite versions).
\end{itemize}

\subsection*{Software Description}
The software implementation section includes:
\begin{itemize}
    \item Installation and configuration of the Arduino IDE for ML model deployment.
    \item Workflow for integrating TensorFlow Lite and TinyML.
    \item Examples of machine learning workflows, including model training, optimization, and deployment on embedded hardware.
\end{itemize}

\subsection*{Development and Testing}
The development process is structured around the Knowledge Discovery in Databases (KDD) framework:
\begin{itemize}
    \item Data selection, preprocessing, and augmentation.
    \item Visualization of the pipeline and validation of machine learning models.
    \item Testing protocols for both hardware and software components, ensuring system robustness and reliability.
\end{itemize}

\subsection*{Deployment and Monitoring}
This section describes the deployment of the gait recognition system, detailing:
\begin{itemize}
    \item The ML pipeline configuration for embedded systems.
    \item Deployment tools, including TensorFlow Lite Micro.
    \item Monitoring and feedback mechanisms to refine the model using real-time data.
\end{itemize}

\subsection*{Concluding Sections}
The report concludes with:
\begin{itemize}
    \item An evaluation of the system's performance and limitations.
    \item Recommendations for future improvements and applications.
    \item References to resources and documentation that support further exploration.
\end{itemize}


\bigskip

